\chapter{Introducción}

\textbf{"Divide y Vencerás"} es un paradigma de programación, que consiste en un conjunto de técnicas algorítmicas en las que primero se divide el problema
a resolver en subproblemas de menor tamaño y de la misma naturaleza, luego se resuelven los subproblemas de manera independiente y recursiva, 
hasta llegar a un caso base (que es un caso tan trivial que detiene la recursión), y finalmente se combinan los resultados para llegar a la solución final. \\

En el presente informe, se analizará empíricamente la complejidad algorítmica de varios algoritmos que utilizan este paradigma de programación, y además,
se evaluará el impacto que tiene cada variante de cada algoritmo. \\

Los algoritmos que serán objeto de este análisis involucrarán tanto algoritmos de ordenamiento, como algoritmos de búsqueda, como también un algoritmo de 
selección. \\

Estos algoritmos son: \\

\textbf{Ordenamiento}

\begin{itemize}
    \item \textit{Merge Sort}.
    \item \textit{Merge-Insertion Sort} (versión de \textit{Merge Sort} que incluye \textit{Insertion Sort} cuando se llegan a tener subarreglos de tamaño pequeño).
    \item \textit{Quick Sort} con pivote en la primera posición.
    \item \textit{Quick Sort} con pivote en la última posición.
    \item \textit{Quick Sort} con pivote en una posición aleatoria.
    \item \textit{Quick Sort} escogiendo al pivote como resultado de una mediana entre 3 elementos.
\end{itemize}

\textbf{Búsqueda}

\begin{itemize}
    \item Búsqueda binaria clásica.
    \item Búsqueda binaria para rangos.
    \item Búsqueda exponencial.
    \item Búsqueda por interpolación.
\end{itemize}

\textbf{Selección}

\begin{itemize}
    \item \textit{Quick Select}.
\end{itemize}

Todo lo anterior se encuentra enmarcado en un sistema de gestión de deportistas desarrollado previamente, el cual trabaja con registros que contienen información 
como ID, nombre, equipo, puntaje y cantidad de competencias disputadas.  Sobre este conjunto de datos, se implementaron distintas funcionalidades prácticas, 
tales como el ordenamiento de deportistas según distintos criterios, la búsqueda eficiente de registros, la obtención del k-ésimo mejor deportista mediante 
algoritmos de selección, y la generación de rankings basados en puntaje.  Además, se realizaron \textit{benchmarks} experimentales utilizando distintos tamaños 
de entrada, con el objetivo de comparar el comportamiento y el rendimiento de los algoritmos implementados en escenarios de mejor caso, peor caso y caso promedio.







